\chapter{सूर्य--पृथ्वी--चंद्रमा तंत्र}\label{ch:g6:sun-earth-moon}

हमारा आकाश तीन दुनियाएँ चलाती हैं: दिन देने वाला सूर्य, हमें ढोने वाली
पृथ्वी, और साथ निभाने वाला \omterm{def:g2:moon-in-the-sky:moon}{चंद्रमा}। इन तीनों को तुम अलग-अलग जानते हो;
अब हम पूरा \emph{तंत्र} बनाएँगे --- कौन किसका चक्कर लगाता है, कितनी
तेज़, कितनी दूर --- और एक पुराना उधार चुकाएँगे: \omterm{def:g3:sun-days-seasons:year}{ऋतुओं} का वादा किया गया
राज़ इसी अध्याय में है।

\section{तीन गतियाँ, तीन घड़ियाँ}

\begin{definition}[अक्ष और कक्षा]\label{def:g6:sun-earth-moon:axisorbit}
पृथ्वी का \emph{अक्ष}\index{अक्ष} उसके ध्रुवों से गुज़रती वह काल्पनिक
रेखा है जिसके चारों ओर वह घूमती है, ठीक लट्टू की कील की तरह।
\emph{कक्षा}\index{कक्षा} वह बंद राह है जिस पर कोई खगोलीय पिंड किसी
दूसरे के चारों ओर चलता है --- सूर्य के चारों ओर पृथ्वी के लिए, और
पृथ्वी के चारों ओर \omterm{def:g2:moon-in-the-sky:moon}{चंद्रमा} के लिए, यह लगभग वृत्त ही है।
\end{definition}

\begin{proposition}[तीन गतियाँ]\label{prop:g6:sun-earth-moon:motions}
पूरा पंचांग तीन सधी हुई गतियों पर टँगा है:
\begin{enumerate}
  \item पृथ्वी अपने \omterm{def:g6:sun-earth-moon:axisorbit}{अक्ष} पर \emph{घूमती} है: $24$ घंटे में एक चक्कर
        --- यही दिन है, जो सूर्य और तारों का जुलूस हमारे आकाश पर
        निकालता है;
  \item पृथ्वी सूर्य के चारों ओर \emph{\omterm{def:g6:sun-earth-moon:axisorbit}{कक्षा} में चलती} है: लगभग
        $365$ दिन और एक चौथाई में एक चक्कर --- यही \omterm{def:g3:sun-days-seasons:year}{वर्ष} है, जो
        सूर्य की चाप को \omterm{def:g3:sun-days-seasons:year}{ऋतुओं} में ऊपर-नीचे सरकाता है;
  \item \omterm{def:g2:moon-in-the-sky:moon}{चंद्रमा} पृथ्वी के चारों ओर \emph{\omterm{def:g6:sun-earth-moon:axisorbit}{कक्षा} में चलता} है: लगभग
        एक महीने में एक चक्कर --- और जैसे-जैसे उसका धूप वाला आधा
        हिस्सा बारी-बारी हमारे सामने आता है, वैसे-वैसे उसकी कलाओं का
        वही जुलूस निकलता है जो तुम्हारी पुरानी चंद्रमा-डायरी में दर्ज
        है।
\end{enumerate}
तीन गतियाँ, तीन घड़ियाँ: दिन, महीना, \omterm{def:g3:sun-days-seasons:year}{वर्ष} --- मानव जाति के सबसे पुराने
समय-मापक, और तीनों आज भी चालू।
\end{proposition}

\begin{omfigure}
\begin{tikzpicture}[scale=0.85]
  \fill[omThm] (5.6,2.2) circle (0.55);
  \foreach \a in {0,30,...,330} {
    \draw[omThm] (5.6,2.2) ++(\a:0.72) -- ++(\a:0.24);
  }
  \node[below, font=\small] at (5.6,1.4) {सूर्य};
  \draw[dashed, omExo, thick] (5.6,2.2) ellipse (4.6 and 1.75);
  \fill[omDef] (10.2,2.2) circle (0.24);
  \draw[-Stealth, thick] (10.2,2.75) arc (90:210:0.5);
  \node[above, font=\small] at (9.7,3.2) {घूर्णन: $24$ घंटे};
  \draw[dashed, omExo] (10.2,2.2) circle (0.85);
  \fill[omExo] (10.8,2.8) circle (0.1);
  \node[right, font=\small] at (10.98,2.8) {चंद्रमा: एक महीना};
  \draw[-Stealth, thick, omDef] (8.1,0.62) arc (222:262:3.9);
  \node[below, font=\small, omDef] at (8.3,0.35) {कक्षा: एक वर्ष};
\end{tikzpicture}

{\small तंत्र का नक़्शा (दूरियाँ हमेशा की तरह दबाई हुई): पृथ्वी सूर्य
के चारों ओर चलते हुए घूमती भी रहती है; और \omterm{def:g2:moon-in-the-sky:moon}{चंद्रमा} पृथ्वी के चारों ओर
\omterm{def:g6:sun-earth-moon:axisorbit}{कक्षा} में चलता है।}
\end{omfigure}

\begin{example}[चौथाई दिन और अधिवर्ष]\label{ex:g6:sun-earth-moon:leap}
पृथ्वी के एक चक्कर में $365$ दिन \emph{और लगभग एक चौथाई} लगता है ---
प्रकृति गोल संख्याओं की वफ़ादार नहीं है। पंचांग यह उधार चौथाइयाँ बचाकर
चुकाते हैं: चार साल की चौथाइयाँ मिलकर एक पूरा दिन बनाती हैं, जो हर
चौथे साल पंचांग में जोड़ दिया जाता है --- यानी $366$ दिन वाला
\emph{अधिवर्ष}, जिसका अतिरिक्त दिन फ़रवरी के अंत में आता है। यह
हिसाब सौ साल तक छोड़ दो और पंचांग \omterm{def:g3:sun-days-seasons:year}{ऋतुओं} से हफ़्तों दूर खिसक जाता है ---
जैसा पुराने पंचांगों ने, तकलीफ़ के साथ, जाना।
\end{example}

\begin{example}[तंत्र, मापनी के साथ]\label{ex:g6:sun-earth-moon:scale}
कॉपी में लिखने लायक संख्याएँ: पृथ्वी लगभग \qty{12800}{km} चौड़ी है;
\omterm{def:g2:moon-in-the-sky:moon}{चंद्रमा} उसका लगभग चौथाई, और वह लगभग \qty{384000}{km} दूर \omterm{def:g6:sun-earth-moon:axisorbit}{कक्षा} में
चलता है --- यानी अगल-बग़ल रखी तीस पृथ्वियों जितनी दूर। सूर्य लगभग
$109$ पृथ्वियों जितना चौड़ा है और $15$ करोड़ किलोमीटर दूर बैठा है। पृथ्वी को सिकोड़कर एक \omterm{def:g2:measuring-length:units}{सेंटीमीटर} की काँच की
गोली बना दो, तो नमूना कहता है: \omterm{def:g2:moon-in-the-sky:moon}{चंद्रमा}, $30$ \omterm{def:g2:measuring-length:units}{सेंटीमीटर} दूर काली मिर्च
का दाना; और सूर्य, एक \omterm{def:g2:measuring-length:units}{मीटर} से चौड़ी गेंद, जो सड़क पर सौ \omterm{def:g2:measuring-length:units}{मीटर} से भी
दूर खड़ी है। हमारा पूरा मुहल्ला ज़्यादातर ख़ालीपन है --- जैसा \omterm{def:g4:solar-system:family}{सौर मंडल}
वाले बाग़ के नमूने ने पहले ही धीरे से कह दिया था।
\end{example}

\section{उधार चुकता: ऋतुएँ क्यों}

\begin{proposition}[झुकाव ही ऋतुएँ बनाता है]\label{prop:g6:sun-earth-moon:tilt}
पृथ्वी का \omterm{def:g6:sun-earth-moon:axisorbit}{अक्ष} उसकी \omterm{def:g6:sun-earth-moon:axisorbit}{कक्षा} पर सीधा खड़ा नहीं है: वह लगभग समकोण के
चौथाई जितना झुका है --- और \omterm{def:g3:sun-days-seasons:year}{वर्ष} के पूरे चक्कर में वह \emph{एक ही ओर}
झुका रहता है, तारों के किसी वफ़ादार \omterm{def:g4:magnets-and-compass:compass}{दिक्सूचक} की तरह। इसलिए चक्कर के
आधे हिस्से में दुनिया का हमारा आधा भाग सूर्य की \emph{ओर} झुका रहता
है: सूर्य ऊँचा चलता है, दिन लंबे होते हैं --- ग्रीष्म। आधा चक्कर बाद
वही न बदला हुआ झुकाव हमें सूर्य से \emph{दूर} कर देता है: नीचा सूर्य,
छोटे दिन --- शिशिर। ऋतुएँ न सूर्य के बदलने से बनती हैं, न दूरी से: वे
एक झुके हुए \omterm{def:g4:solar-system:planet}{ग्रह} की देन हैं, जो अपना न बदला हुआ झुकाव लिए अपने तारे
के चारों ओर घूमता रहता है।
\end{proposition}

\begin{omfigure}
\begin{tikzpicture}[scale=0.85]
  \fill[omThm] (5.6,2.0) circle (0.5);
  \foreach \a in {0,45,...,315} {
    \draw[omThm] (5.6,2.0) ++(\a:0.65) -- ++(\a:0.22);
  }
  \begin{scope}[shift={(1.1,2.0)}]
    \draw[thick] (0,0) circle (0.62);
    \draw[very thick, omDef] (-0.28,-0.75) -- (0.28,0.75);
    \fill[omDef] (0.19,0.5) circle (0.09);
  \end{scope}
  \node[below, font=\small, align=center] at (1.1,0.9)
    {सूर्य की ओर झुका:\\यहाँ ग्रीष्म};
  \begin{scope}[shift={(10.1,2.0)}]
    \draw[thick] (0,0) circle (0.62);
    \draw[very thick, omDef] (-0.28,-0.75) -- (0.28,0.75);
    \fill[omDef] (0.19,0.5) circle (0.09);
  \end{scope}
  \node[below, font=\small, align=center] at (10.1,0.9)
    {आधा चक्कर बाद, वही झुकाव:\\दूर की ओर --- यहाँ शिशिर};
\end{tikzpicture}

{\small एक ही झुकी हुई पृथ्वी, अपने चक्कर के दो पड़ावों पर। \omterm{def:g6:sun-earth-moon:axisorbit}{अक्ष}
(नीला) दोनों में एक ही ओर झुका है --- चिह्नित उत्तरी क़स्बे के लिए जून
में सूर्य की ओर, और दिसंबर में उससे दूर।}
\end{omfigure}

\begin{method}[गोला और लैंप]\label{met:g6:sun-earth-moon:globe}
पूरी मशीन मेज़ पर देखो:
\begin{enumerate}
  \item मेज़ के बीच सूर्य की जगह एक लैंप रखो, और एक झुका हुआ गोला
        (या तिरछी सींक पर टिकी गेंद) उसके चारों ओर धीरे-धीरे ले
        चलो, \emph{और पूरे रास्ते सींक को कमरे के एक ही कोने की ओर
        ताके रखो};
  \item चक्कर के हर चौथाई पर रुको और अपने घर वाले अक्षांश को देखो:
        उसके रोज़ के घूर्णन-वृत्त का कितना हिस्सा उजाले में पड़ता
        है;
  \item लैंप की एक ओर तुम्हारा गोलार्ध झुककर धूप सेंकता है ---
        उजाले वाली लंबी चापें, ऊँचा लैंप: ग्रीष्म; और सामने वाली ओर,
        दूर झुका हुआ --- छोटी चापें, नीचा लैंप: शिशिर;
  \item बीच के दोनों पड़ावों पर कोई भी ध्रुव भीतर की ओर नहीं झुकता:
        दिन और रात बराबर बँट जाते हैं --- वसंत और शरद।
\end{enumerate}
इसे चलाने के लिए एक ही अनुशासन चाहिए: सींक को कभी घूमने मत दो ---
\omterm{def:g6:sun-earth-moon:axisorbit}{अक्ष} की वही वफ़ादारी \emph{ही} ऋतुएँ है।
\end{method}

\begin{example}[पुराने सुराग़ों का हिसाब]\label{ex:g6:sun-earth-moon:clues}
\omterm{def:g3:sun-days-seasons:sundial}{धूपघड़ी} वाले वर्षों का तुम्हारा हर \omterm{def:g1:five-senses:observation}{प्रेक्षण} अब झुकाव को रिपोर्ट करता
है। ग्रीष्म की ऊँची चाप: हमारा गोलार्ध भीतर की ओर झुका है, इसलिए
दोपहर में सूर्य ज़्यादा ऊँचा खड़ा होता है। शिशिर की ठिगनी चाप और लंबी
छायाएँ: दूर की ओर झुके होने पर सूर्य नीचे-नीचे सरक जाता है। भूमध्य
रेखा के उस पार उलटी ऋतुएँ --- दक्षिण में दिसंबर की समुद्र-तट वाली
छुट्टियाँ जबकि उत्तर हिम हटाता है: जब एक गोलार्ध भीतर की ओर झुकता है,
तब दूसरे को दूर झुकना ही पड़ता है। वसंत और शरद के बराबर दिन भी लैंप और
गोले वाली सैर से अपने आप निकल आते हैं। समकोण के एक चौथाई जितना अकेला
झुकाव पूरा तमाशा चलाता है।
\end{example}

\begin{remark}[दूरी नहीं --- और उसका सबूत]\label{rem:g6:sun-earth-moon:notdistance}
पृथ्वी की \omterm{def:g6:sun-earth-moon:axisorbit}{कक्षा} इतनी वृत्त जैसी है कि सूर्य से दूरी बहुत कम बदलती है
--- और मज़े की बात यह कि उत्तरी शिशिर में हम सूर्य के ज़रा
\emph{ज़्यादा पास} होते हैं। \omterm{def:g3:sun-days-seasons:year}{ऋतुओं} को यात्रा की लंबाई से नहीं, \omterm{def:g1:light-and-shadows:source}{प्रकाश}
के \emph{तिरछेपन} से मतलब है: भीतर की ओर झुके हों तो सूर्य का \omterm{def:g1:light-and-shadows:source}{प्रकाश}
खड़ा और सिमटा हुआ पड़ता है; दूर झुके हों तो वही \omterm{def:g1:light-and-shadows:source}{प्रकाश} ज़मीन पर पतला
फैल जाता है। दीवार पर टॉर्च सीधी और फिर तिरछी करके देखो, फ़र्क़ एक
सेकंड में दिख जाएगा --- खड़ा \omterm{def:g1:light-and-shadows:source}{प्रकाश} काटता है, ढलवाँ \omterm{def:g1:light-and-shadows:source}{प्रकाश} बमुश्किल
गरम करता है।
\end{remark}

\section{अभ्यास}

\begin{exercise}[$\star$]\label{exo:g6:sun-earth-moon:1}
तंत्र की तीनों गतियाँ बताओ और यह भी कि हर एक कौन-सी घड़ी चलाती है।
\end{exercise}

\begin{exercise}[$\star$]\label{exo:g6:sun-earth-moon:2}
पृथ्वी का \omterm{def:g6:sun-earth-moon:axisorbit}{अक्ष} क्या है? उसके बारे में समकोण के चौथाई वाला कौन-सा तथ्य,
और साथ में कौन-सी वफ़ादारी, ऋतुएँ बना देते हैं?
\end{exercise}

\begin{exercise}[$\star$]\label{exo:g6:sun-earth-moon:3}
पंचांग हर चौथे साल एक दिन क्यों जोड़ता है? अगर वह कभी न जोड़ता तो क्या
खिसक जाता?
\end{exercise}

\begin{exercise}[$\star$]\label{exo:g6:sun-earth-moon:4}
अगल-बग़ल रखी लगभग कितनी पृथ्वियाँ \omterm{def:g2:moon-in-the-sky:moon}{चंद्रमा} तक पहुँचती हैं? और सूर्य के
चेहरे पर कितनी पृथ्वियाँ समाती हैं?
\end{exercise}

\begin{exercise}[$\star$]\label{exo:g6:sun-earth-moon:5}
गोला-और-लैंप वाली सैर में सींक को कमरे के एक ही कोने की ओर ताके रखना
क्यों ज़रूरी है? वह घूम जाए तो क्या टूट जाता है?
\end{exercise}

\begin{exercise}[$\star$]\label{exo:g6:sun-earth-moon:6}
जब यहाँ ग्रीष्म होता है, तब भूमध्य रेखा के उस पार वाले गोलार्ध में
कौन-सी \omterm{def:g3:sun-days-seasons:year}{ऋतु} होती है, और ऐसा होना ज़रूरी क्यों है?
\end{exercise}

\begin{exercise}[$\star$]\label{exo:g6:sun-earth-moon:7}
\cref{rem:g6:sun-earth-moon:notdistance} वाली दीवार-और-टॉर्च की
तसवीर से समझाओ कि सूर्य की ओर झुकने से हमें ज़्यादा गरमी क्यों मिलती
है।
\end{exercise}

\begin{exercise}[$\star$]\label{exo:g6:sun-earth-moon:8}
एक पुराना दोस्त लौटता है: महीने भर में \omterm{def:g2:moon-in-the-sky:moon}{चंद्रमा} अपनी कलाएँ क्यों
दिखाता है? तीनों गतियों में से कौन-सी यह तमाशा चलाती है?
\end{exercise}

\begin{exercise}[$\star\star$]\label{exo:g6:sun-earth-moon:9}
एक ज़िद्दी चाचा अड़े हैं: ``शिशिर तब होता है जब पृथ्वी सूर्य से सबसे
दूर होती है।'' इस अध्याय के ऐसे दो तथ्य जुटाओ जो मिलकर उन्हें हरा
दें।
\end{exercise}

\begin{exercise}[$\star\star$]\label{exo:g6:sun-earth-moon:10}
मान लो पृथ्वी का \omterm{def:g6:sun-earth-moon:axisorbit}{अक्ष} उसकी \omterm{def:g6:sun-earth-moon:axisorbit}{कक्षा} पर बिलकुल सीधा खड़ा होता --- ज़रा भी
झुकाव नहीं। एक \omterm{def:g3:sun-days-seasons:year}{वर्ष} का वर्णन करो: \omterm{def:g3:sun-days-seasons:year}{ऋतुओं} का क्या होता, दोपहर में सूर्य
की ऊँचाई का क्या, और दिनों की लंबाई का क्या? (मन ही मन सीधे खड़े गोले
को लैंप के चारों ओर घुमाओ।)
\end{exercise}

\begin{exercise}[$\star\star$]\label{exo:g6:sun-earth-moon:11}
\omterm{def:g2:moon-in-the-sky:moon}{चंद्रमा} का एक चक्कर लगभग $27$ दिन का होता है, फिर भी एक \omterm{ex:g2:moon-in-the-sky:shapes}{पूर्ण चंद्र} से
अगले \omterm{ex:g2:moon-in-the-sky:shapes}{पूर्ण चंद्र} तक लगभग $29.5$ दिन लगते हैं। इसका हल सुझाओ ---
\omterm{def:g2:moon-in-the-sky:moon}{चंद्रमा} के चक्कर लगाने के दौरान और क्या खिसक गया? (एक चंद्र-चक्कर के
दौरान अपनी \omterm{def:g6:sun-earth-moon:axisorbit}{कक्षा} पर आगे बढ़ती पृथ्वी का रेखाचित्र बनाओ।)
\end{exercise}

\begin{exercise}[$\star\star\star$]\label{exo:g6:sun-earth-moon:12}
भरी गरमी में \omterm{def:g4:magnets-and-compass:poles}{उत्तरी ध्रुव} के पास के इलाक़ों में आधी रात को भी सूर्य
क्षितिज के ऊपर बना रहता है --- वही मशहूर ``आधी रात का सूर्य'' --- और
भरे जाड़े में वे हफ़्तों तक किसी \omterm{def:g1:day-and-night:daynight}{सूर्योदय} का इंतज़ार करते हैं।
\cref{met:g6:sun-earth-moon:globe} वाले झुके घूर्णन-वृत्तों से दोनों
अचरज समझाओ: जून में और दिसंबर में झुकाव किसी ध्रुवीय क़स्बे के रोज़ के
वृत्त के साथ क्या करता है?
\end{exercise}

\section{समस्या: खेल के मैदान वाला नमूना}

\begin{problem}[{सप्ताहांत समस्या --- कक्षा खेल के मैदान पर
सूर्य--पृथ्वी--चंद्रमा तंत्र बनाती है; काँच की गोलियाँ, काली मिर्च के
दाने और एक बहुत लंबी सैर; और नमूने में शामिल होता झुकाव}]\label{pb:g6:sun-earth-moon:1}
\omterm{def:g6:sun-earth-moon:axisorbit}{कक्षा} तंत्र को मापनी के अनुसार बनाती है: पृथ्वी \qty{1}{cm} चौड़ी काँच
की गोली है। काम की सच्ची संख्याएँ: पृथ्वी \qty{12800}{km} चौड़ी;
\omterm{def:g2:moon-in-the-sky:moon}{चंद्रमा} \qty{3200}{km} चौड़ा, और \qty{384000}{km} दूर \omterm{def:g6:sun-earth-moon:axisorbit}{कक्षा} में; सूर्य
$109$ पृथ्वियों जितना चौड़ा, और $15$ करोड़ किलोमीटर दूर।

\medskip\noindent\textbf{भाग I --- भूमिकाएँ बाँटना।}
\begin{enumerate}
  \item नमूने में \qty{1}{cm} कितने किलोमीटर का प्रतिनिधि है?
  \item \omterm{def:g2:moon-in-the-sky:moon}{चंद्रमा} पृथ्वी की चौड़ाई का लगभग चौथाई है। नमूने वाला \omterm{def:g2:moon-in-the-sky:moon}{चंद्रमा}
        कितना चौड़ा हुआ --- और रसोई की कौन-सी चीज़ उसकी भूमिका अच्छी
        निभाएगी?
  \item \omterm{def:g2:moon-in-the-sky:moon}{चंद्रमा} की सच्ची दूरी $30$ पृथ्वी-चौड़ाइयों की है। नमूने वाला
        \omterm{def:g2:moon-in-the-sky:moon}{चंद्रमा} गोली से कहाँ खड़ा होगा?
  \item नमूने वाला सूर्य $109$ गोलियों जितना चौड़ा है। यह कितने
        \omterm{def:g2:measuring-length:units}{सेंटीमीटर} हुआ --- और मैदान की कौन-सी चीज़ लगभग इतनी बड़ी
        होती है?
\end{enumerate}

\medskip\noindent\textbf{भाग II --- लंबी सैर।}
\begin{enumerate}[resume]
  \item सूर्य की दूरी $15$ करोड़ किलोमीटर है। नमूने में --- पहले सवाल
        के उत्तर से भाग देकर --- यह कितने \omterm{def:g2:measuring-length:units}{सेंटीमीटर} हुई, और कितने
        \omterm{def:g2:measuring-length:units}{मीटर}?
  \item क्या यह नमूना \qty{100}{m} के खेल मैदान पर समा जाएगा?
        सूर्य वाली गेंद कहाँ खड़ी करनी होगी?
  \item एक विद्यार्थी पृथ्वी--सूर्य वाला हिस्सा $1$ \omterm{def:g2:measuring-length:units}{मीटर} प्रति सेकंड
        की \omterm{def:g5:speed-distance-time:speed}{चाल} से चलता है। इस सैर में लगभग कितने मिनट लगेंगे?
  \item गोली के पास खड़े होकर $117$ \omterm{def:g2:measuring-length:units}{मीटर} दूर रखी एक \omterm{def:g2:measuring-length:units}{मीटर} चौड़ी सूर्य
        वाली गेंद की ओर देखते हुए \omterm{def:g6:sun-earth-moon:axisorbit}{कक्षा} चौंक जाती है: वह लगभग उतनी ही
        बड़ी दिखती है जितना गोली से $30$ \omterm{def:g2:measuring-length:units}{सेंटीमीटर} दूर रखा चौथाई
        \omterm{def:g2:measuring-length:units}{सेंटीमीटर} वाला चंद्रमा-दाना। दोनों अनुपात जाँचो --- और इस
        संयोग को उस चीज़ से जोड़ो जो असली आकाश में कभी-कभार, और बड़ी
        शान से, होती है।
\end{enumerate}

\medskip\noindent\textbf{भाग III --- नमूने को चलाना।}
\begin{enumerate}[resume]
  \item नमूने को ईमानदारी से चलाना हो और एक सच्चा दिन एक सेकंड में
        निभाया जाए, तो गोली को सूर्य वाली गेंद का चक्कर लगाने में
        कितना समय लगना चाहिए, और मटर को गोली का चक्कर लगाने में?
        (मोटे तौर पर गिनो: \omterm{def:g3:sun-days-seasons:year}{वर्ष} $365$ दिन, चंद्र-चक्कर $27$ दिन।)
  \item गोली को घूमना भी है। एक सच्चा दिन प्रति सेकंड की दर पर नमूने
        में वह कितनी तेज़ हुई --- और \qty{1}{cm} की गोली पर उसे कोई
        क्यों नहीं देख पाएगा?
  \item गोली के भीतर से गुज़री एक सींक \omterm{def:g6:sun-earth-moon:axisorbit}{अक्ष} की भूमिका निभाती है। पूरे
        चक्कर में उस सींक को किस झुकाव पर और किस ओर ताके हुए ले जाना
        होगा?
  \item नमूने की ईमानदार चेतावनी लिखो: एक बात बताओ जो खेल-मैदान वाला
        नमूना सच-सच दिखा देता है, और एक बात जो मैदान भर का कोई नमूना
        एक साथ नहीं दिखा सकता (आकार और दूरियाँ साथ-साथ? गतियाँ?
        \omterm{def:g1:light-and-shadows:source}{प्रकाश}?)। अपनी संख्याओं से तर्क दो।
\end{enumerate}
\end{problem}
